% !TeX program = xelatex
\documentclass[12pt,a4paper]{article}
\usepackage{fontspec}
\setmainfont{Times New Roman}
\usepackage[english,russian]{babel}
\usepackage{amsmath,amssymb,mathrsfs,cancel}
\usepackage[margin=22mm]{geometry}
\usepackage{graphicx,float}
\usepackage[unicode,hidelinks]{hyperref}
\DeclareMathOperator{\grad}{grad}
\DeclareMathOperator{\Div}{div}
\newcommand{\Rey}{\mathrm{Re}}
\setlength{\emergencystretch}{3em}
\sloppy
\interfootnotelinepenalty=10000
\newcommand{\topic}[1]{\par\medskip\noindent\textbf{#1}\par\nobreak\smallskip}
\renewcommand{\thefootnote}{\arabic{footnote})}
\begin{document}
\section*{Вопрос 12}
Понятие об эволюционности разрывов для гиперболических уравнений. Эволюционность скачка насыщенности в задаче Бакли — Леверетта. Модель двухфазных течений с учетом капиллярной разности давлений. Понятие об уравнении Рапопорта — Лиса. Представление о задаче о структуре скачка насыщенности.

\medskip
\noindent\textbf{Источник:} А. Г. Куликовский, Е. И. Свешникова. \textit{Нелинейные волны в упругих телах}. Москва: Московский лицей, 1998. § 1.3, с. 26--31; К. С. Басниев, А. М. Власов, И. Н. Кочина, В. М. Максимов. \textit{Подземная гидравлика}. Москва: Недра, 1986. гл. 9, фрагменты с. 207, 210, 219--220; Г. И. Баренблатт, В. М. Ентов, В. М. Рыжик. \textit{Теория нестационарной фильтрации жидкости и газа}. Москва: Недра, 1972. гл. VI, § 2 (определения), § 3 (начало), с. 157, 159, 164--167; рисунок VI.11 на с. 168.
\par\bigskip
\begin{quote}\sloppy
Закон движения границы $x=X(t)$ на плоскости $x,t$ представлен жирной линией на рис. 1.2, наклон которой определен скоростью $W$. Будем сначала предполагать, что скорость границы $W$ не совпадает ни с одной из характеристических скоростей системы (1.6), т.е. $W\ne c^{(m)}$, $m=1,\ldots,n$. Тогда все $n$ характеристик (и соответствующие им бегущие волны) в каждой точке границы можно разделить на \textit{характеристики, уходящие от границы} в область $x>X(t)$, для которых $c^{(r)}>W$, $r=1,2,\ldots,s$, и \textit{характеристики, приходящие к границе}, для которых $c^{(l)}<W$, $l=s+1,\ldots,n$. На рис. 1.2 характеристики изображены прямыми со стрелками в направлении роста $t$. На приходящих характеристиках известны из начальных условий значения инвариантов Римана $w_l$. Для построения решения (1.12) нужно знать значения остальных функций $w_r$ ($r=1,\ldots,s$) на уходящих от границы характеристиках.

Таким образом, для однозначного продолжения решения в области $t>t_0$ нужно так задать граничные условия, чтобы на кривой $x=X(t)$ они позволили найти $s$ значений инвариантов Римана уходящих волн $w_r$ по известным значениям инвариантов Римана приходящих волн $w_l$, т.е. $w_r=\varphi_r(w_{s+1},\ldots,w_n)$, $r=1,2,\ldots,s$. Для этого число соотношений, задающих связи между функциями $w_j$ плюс закон движения границы должно быть равно $s+1$.

Часто в задачах механики движение границы не задано, но имеется выражение для ее скорости $W$ через функции $u_i$, а связи между неизвестными функциями на границе включают и величину $W$. В этом случае все $s+1$ соотношения становятся равноправными. В исходных переменных $u_i$ они могут быть записаны в виде
\begin{equation}\Phi_q(u_1,\ldots,u_k,W)=0.\tag{1.16}\label{q12:eq-1-16}\end{equation}


Эти соотношения должны позволить найти $s+1$ неизвестную величину $w_1,w_2,\ldots,w_s,W$. Кроме того, следствием этих уравнений не должно быть никакое соотношение, связывающее между собой амплитуды приходящих волн $w_{s+1},\ldots,w_n$, так как эта связь будет противоречить произвольности начальных условий, определяющих эти величины. Чтобы это было возможно, необходимо, чтобы система (1.16) содержала $s+1$ уравнение с отличным от нуля якобианом относительно $w_r$ и $W$, $r=1,2,\ldots,s$.

Таким образом, \textit{для однозначной разрешимости начально-краевой задачи для линейной гиперболической системы необходимо задавать $s+1$ граничных условий вида} (1.16), обладающих отличным от нуля якобианом относительно $w_1,w_2,\ldots,w_s$ и $W$, где $s$ --- число характеристик, уходящих от границы в рассматриваемой точке $x,t$. Такие \textit{граничные условия называются эволюционными}. Часто говорят об эволюционной границе, имея в виду то же самое.
\par\nobreak\smallskip{\small А. Г. Куликовский, Е. И. Свешникова. \textit{Нелинейные волны в упругих телах}. Москва: Московский лицей, 1998. Глава 1, § 1.3, с. 27--28.}\par
\end{quote}

\begin{quote}\sloppy
В некоторых случаях граничные условия должны выставляться на границе, разделяющей области, в каждой из которых решение описывается своей гиперболической системой (это могут быть, в частности, системы и с различным числом неизвестных). В этом случае условие эволюционности границы остается прежним, однако под $s$ следует понимать число всех уходящих в обе стороны от границы характеристик (налево уходят характеристики, у которых скорость меньше, чем $W$).
\par\nobreak\smallskip{\small А. Г. Куликовский, Е. И. Свешникова. \textit{Нелинейные волны в упругих телах}. Москва: Московский лицей, 1998. Глава 1, § 1.3, с. 29.}\par
\end{quote}

\begin{quote}\sloppy
\textit{Для нелинейных задач под условиями эволюционности границы подразумевается выполнение условий эволюционности для линеаризованной задачи}. Линеаризация производится в малой окрестности точки на границе, в которой $u_i=u_i^0=\mathrm{const}$, так что $u_i=u_i^0+v_i(x,t)$, при этом $v_i$ малы.

Эволюционность означает существование однозначного решения задачи о взаимодействии границы с малыми возмущениями в рамках линеаризованной постановки задачи для функций $v_i$ и $\delta W$, где $\delta W$ --- малое возмущение скорости границы.
\par\nobreak\smallskip{\small А. Г. Куликовский, Е. И. Свешникова. \textit{Нелинейные волны в упругих телах}. Москва: Московский лицей, 1998. Глава 1, § 1.3, с. 30.}\par
\end{quote}

\begin{quote}\sloppy
До сих пор предполагалось, что скорость границы $W$ не равна ни одной из характеристических скоростей. Если такое равенство имеет место, то будем говорить, что на границе выполнено \textit{условие Жуге} (по аналогии с теорией детонационных волн (Ландау и Лившиц [1986])). К такого типа границам относятся некоторые разрывы в задачах механики сплошных сред, играющие важную роль при построении решений. О них будет подробно сказано при описании решений с разрывами в упругих средах (Главы 4 и 5).
\par\nobreak\smallskip{\small А. Г. Куликовский, Е. И. Свешникова. \textit{Нелинейные волны в упругих телах}. Москва: Московский лицей, 1998. Глава 1, § 1.3, с. 31.}\par
\end{quote}

\begin{quote}\sloppy
\begin{equation}f(\sigma)=\frac{k_1(\sigma)}{k_1(\sigma)+\mu_0k_2(\sigma)},\qquad\mu_0=\mu_1/\mu_2.\tag{9.22}\label{q12:eq-9-22}\end{equation}

\par\nobreak\smallskip{\small К. С. Басниев, А. М. Власов, И. Н. Кочина, В. М. Максимов. \textit{Подземная гидравлика}. Москва: Недра, 1986. Глава 9, § 4, с. 207.}\par
\end{quote}

\begin{quote}\sloppy
В случае одномерного течения несжимаемых несмешивающихся жидкостей в условиях, когда можно пренебречь капиллярным давлением, а также влиянием силы тяжести, процесс вытеснения допускает простое математическое описание.

Для обоих случаев одномерного потока (прямолинейно-параллельного и плоскорадиального) это приводит к классической в теории вытеснения модели Баклея--Леверетта, описываемой однотипным уравнением для насыщенности $\sigma$ вытесняющей фазы 1, которое получается из (9.34) при гравитационном параметре $A_n=0$ и имеет вид
\begin{equation}\frac{\partial\sigma}{\partial\tau}+f'(\sigma)\frac{\partial\sigma}{\partial\xi}=0.\tag{9.35}\label{q12:eq-9-35}\end{equation}
\par\nobreak\smallskip{\small К. С. Басниев, А. М. Власов, И. Н. Кочина, В. М. Максимов. \textit{Подземная гидравлика}. Москва: Недра, 1986. Глава 9, § 5, с. 210.}\par
\end{quote}

\begin{quote}\sloppy
Вообще говоря, к перечисленным условиям необходимо добавить еще одно --- определяющее так называемую устойчивость разрывного решения. Для устойчивости скачка необходимо, чтобы в любой точке кривой $\xi_c(\tau)$ пересекались две характеристики уравнения (9.35). При этом любая характеристика, скорость которой $f'(\sigma)$ равна скорости скачка $D$, считается приходящей на скачок. Это условие для класса рассматриваемых задач с кусочно-постоянными данными (9.39) можно сформулировать следующим образом.

Пусть $L$ --- прямая на плоскости $(\sigma,f(\sigma))$, соединяющая точки $(\sigma^-,f^-)$ и $(\sigma^+,f^+)$. Здесь, как и ранее, $\sigma^-$ --- значение насыщенности «до» разрыва (правосторонний предел по отношению к изменению $\xi$), а $\sigma^+$ --- значение насыщенности «после» разрыва (левосторонний предел); $f^-=f(\sigma^-)$ и $f^+=f(\sigma^+)$ соответственно (см.
рис. 9.4). Тогда вдоль допустимого разрыва кривая $f(\sigma)$ на интервале $(\sigma^-,\sigma^+)$ располагается ниже прямой $L$, если $\sigma^-<\sigma^+$, и выше прямой $L$, если $\sigma^->\sigma^+$.
\par\nobreak\smallskip{\small К. С. Басниев, А. М. Власов, И. Н. Кочина, В. М. Максимов. \textit{Подземная гидравлика}. Москва: Недра, 1986. Глава 9, § 5, с. 219--220.}\par
\end{quote}

\begin{quote}\sloppy
В систему уравнений и граничные условия входят следующие размерные определяющие параметры: $\Delta p$ --- перепад давления между скважинами или галереями; $L$ --- характерный размер области течения; коэффициенты $k/\mu_1$ и $k/\mu_2$, а также $p_c^0=\alpha\sqrt{m/k}\cos\Theta$. Перейдем в уравнениях (VI.2.1), (VI.2.2), (VI.2.6) к безразмерным переменным:
\[\begin{gathered}
\frac{p_i}{\Delta p}=P_i;\quad\bar U_i=\frac{\bar u_i}{u_0};\quad\left(\bar u_0=\frac{k\Delta p}{\mu_1L}\right);\quad X=\frac{x}{L};\\
Y=\frac{y}{L};\quad Z=\frac{z}{L};\quad\tau=\frac{u_0t}{L}=\frac{k\Delta p}{\mu_1L^2}t;\quad\varepsilon=\frac{p_c^0}{\Delta p}.
\end{gathered}\]

\par\nobreak\smallskip{\small Г. И. Баренблатт, В. М. Ентов, В. М. Рыжик. \textit{Теория нестационарной фильтрации жидкости и газа}. Москва: Недра, 1972. Глава VI, § 2, с. 157.}\par
\end{quote}

\begin{quote}\sloppy
где $q(t)=-\dfrac{k}{\mu_1}C(t)$ --- суммарный расход жидкости через трубку тока, а
\[F(s)=\frac{f_1(s)}{f_1(s)+\mu_0f_2(s)}.\]
Функция $F(s)$ равна отношению скорости фильтрации (или расхода) вытесняющей фазы к суммарной скорости фильтрации (или к суммарному расходу). Функцию $F(s)$ принято называть функцией распределения фаз.
\par\nobreak\smallskip{\small Г. И. Баренблатт, В. М. Ентов, В. М. Рыжик. \textit{Теория нестационарной фильтрации жидкости и газа}. Москва: Недра, 1972. Глава VI, § 2, с. 159.}\par
\end{quote}

\begin{quote}\sloppy
На примере решения Баклея --- Леверетта видно, что если в задаче о вытеснении нефти водой ограничиться первым членом внешнего разложения, т. е. не учитывать капиллярной разности давлений, то необходимо вводить поверхности, на которых насыщенность терпит разрыв.


\par\nobreak\smallskip{\small Г. И. Баренблатт, В. М. Ентов, В. М. Рыжик. \textit{Теория нестационарной фильтрации жидкости и газа}. Москва: Недра, 1972. Глава VI, § 3, с. 164.}\par
\end{quote}

\begin{quote}\sloppy
Чтобы описать распределение насыщенности в узких областях, вблизи скачка, нужно получить «внутреннее разложение» решения задачи о вытеснении на основе полной системы уравнений (VI.2.7), (VI.2.8). Для построения первого члена внутреннего разложения введем в окрестности некоторой точки поверхности разрыва (т. е. скачка насыщенности) локальную мгновенную декартову систему координат с центром в точке $O$ поверхности разрыва. Ось $x$ направим по нормали к поверхности разрыва и введем вдоль этой оси масштаб $l=\varepsilon L$, т. е. положим $X=x/l$, сохраняя вдоль других осей масштаб $L$ (рис. VI.10). Масштаб времени примем равным $t_1=l/u_0$ и положим $\tau=t/t_1$. В остальном сохраним те же безразмерные параметры и переменные, что и в уравнениях (VI.2.7) и (VI.2.8).


Запишем систему уравнений в безразмерном виде. Обобщенный закон Дарси:
\begin{equation}\begin{aligned}
U_{1X}&=-\frac1\varepsilon f_1(s)\frac{\partial P_1}{\partial X},&U_{2X}&=-\frac{\mu_0}{\varepsilon}f_2(s)\frac{\partial P_2}{\partial X};\\
U_{1Y}&=-f_1(s)\frac{\partial P_1}{\partial Y},&U_{2Y}&=-\mu_0f_2(s)\frac{\partial P_2}{\partial Y};\\
U_{1Z}&=-f_1(s)\frac{\partial P_1}{\partial Z},&U_{2Z}&=-\mu_0f_2(s)\frac{\partial P_2}{\partial Z};\\
P_2-P_1&=\varepsilon J(s),&\mu_0&=\frac{\mu_1}{\mu_2}.
\end{aligned}\tag{VI.3.1}\label{q12:eq-VI-3-1}\end{equation}
Уравнения непрерывности:
\begin{equation}\begin{gathered}
m\frac{\partial s}{\partial\tau}+\frac{\partial U_{1X}}{\partial X}+\varepsilon\frac{\partial U_{1Y}}{\partial Y}+\varepsilon\frac{\partial U_{1Z}}{\partial Z}=0;\\
m\frac{\partial s}{\partial\tau}-\frac{\partial U_{2X}}{\partial X}-\varepsilon\frac{\partial U_{2Y}}{\partial Y}-\varepsilon\frac{\partial U_{2Z}}{\partial Z}=0.
\end{gathered}\tag{VI.3.2}\label{q12:eq-VI-3-2}\end{equation}
Подставляя (VI.3.1) в (VI.3.2) и отбрасывая члены порядка $\varepsilon$ и $\varepsilon^2$ получим
\begin{equation}\begin{gathered}
m\frac{\partial s}{\partial\tau}-\frac1\varepsilon\frac{\partial}{\partial X}\left(f_1(s)\frac{\partial P_1}{\partial X}\right)=0;\\
m\frac{\partial s}{\partial\tau}+\frac{\mu_0}{\varepsilon}\frac{\partial}{\partial X}\left(f_2(s)\frac{\partial P_1}{\partial X}\right)+\mu_0\frac{\partial}{\partial X}\left(f_2(s)J'(s)\frac{ds}{dX}\right)=0.
\end{gathered}\tag{VI.3.3}\label{q12:eq-VI-3-3}\end{equation}
Исключим $\partial P_1/\partial X$ из уравнений (VI.3.3), что дает окончательно
\begin{equation}m\frac{\partial s}{\partial\tau}-w\frac{\partial F(s)}{\partial X}+\mu_0\frac{\partial}{\partial X}\left(f_2(s)F(s)J'(s)\frac{\partial s}{\partial X}\right)=0,\tag{VI.3.4}\label{q12:eq-VI-3-4}\end{equation}
где
\begin{equation}\begin{aligned}
w(\tau)&=\frac{u_1+u_2}{u_0}=-\frac1\varepsilon[f_1(s)+\mu_0f_2(s)]\frac{\partial P_1}{\partial X}+\mu_0f_2J'\frac{\partial s}{\partial X}=\\
&=U_{1X}+U_{2X}.
\end{aligned}\tag{VI.3.5}\label{q12:eq-VI-3-5}\end{equation}
Система (VI.3.1)--(VI.3.2) свелась к одномерному уравнению (VI.3.4), потому что радиус кривизны поверхности разрыва имеет порядок $L$, и в принятом масштабе $l$ эта поверхность, как и поверхности $s=\mathrm{const}$, заменяется плоскостями. Равенство (VI.3.5) означает, что в пределах зоны скачка, где движение можно считать одномерным, суммарная скорость фильтрации обеих фаз вдоль оси $x$, $w$ есть величина, зависящая только от времени, как при движении в цилиндрической трубке тока.

Значение $w$ в уравнении (VI.3.4) находится из внешнего разложения, как безразмерная суммарная скорость фильтрации через поверхность разрыва в точке $O$. Уравнение (VI.3.4), описывающее одномерное вытеснение несмешивающихся жидкостей, называется уравнением Рапопорта --- Лиса [152].
\par\nobreak\smallskip{\small Г. И. Баренблатт, В. М. Ентов, В. М. Рыжик. \textit{Теория нестационарной фильтрации жидкости и газа}. Москва: Недра, 1972. Глава VI, § 3, с. 165--166.}\par
\end{quote}

\begin{quote}\sloppy
Поскольку масштаб времени во внутреннем разложении $t_1$ намного меньше, чем во внешнем, а скорость $w$ определяется внешним разложением, то при исследовании внутреннего разложения можно считать $w(\tau)=\mathrm{const}$. Ввиду различия масштабов времени во внешнем и внутреннем разложении достаточно воспользоваться стационарным решением задачи Коши для уравнения (VI.3.4), т. е. положить
\begin{equation}s=s(\bar x),\quad\bar x=X-V^0\tau\quad\left(V^0=\frac{V}{u_0}\right).\tag{VI.3.6}\label{q12:eq-VI-3-6}\end{equation}
Иными словами, в масштабе времени внутреннего разложения процесс вытеснения продолжался весьма долго, и его можно считать установившимся в системе координат, связанной со скачком. При этом в силу различия масштабов $l$ и $L$ должны выполняться граничные условия
\begin{equation}s(-\infty)=s^{(1)}=s_c;\quad s(+\infty)=s^{(2)}=s_0,\tag{VI.3.7}\label{q12:eq-VI-3-7}\end{equation}
где $s^{(1)}=s_c$ и $s^{(2)}=s_0$ --- насыщенности за и перед скачком, определяемые из внешнего разложения [в задаче Баклея --- Леверетта они связаны соотношением (VI.2.30)]. Параметр $V^0$ в равенстве (VI.3.6) есть, очевидно, скорость распространения поверхности разрыва, определяемая формулами (VI.2.21) или (VI.2.24).

Используя (VI.3.6), получим вместо (VI.3.4) уравнение
\begin{equation}-mV^0\frac{ds}{d\bar x}-w\frac{dF(s)}{d\bar x}+\mu_0\frac{d}{d\bar x}\left[f_2(s)F(s)J'(s)\frac{ds}{d\bar x}\right]=0.\tag{VI.3.8}\label{q12:eq-VI-3-8}\end{equation}

\par\nobreak\smallskip{\small Г. И. Баренблатт, В. М. Ентов, В. М. Рыжик. \textit{Теория нестационарной фильтрации жидкости и газа}. Москва: Недра, 1972. Глава VI, § 3, с. 166.}\par
\end{quote}

\end{document}
